\documentclass[11pt]{article}
\usepackage[a4paper,margin=2.35cm]{geometry}
\usepackage{fontspec}
\usepackage{amsmath,amssymb}
\setmainfont[ItalicFont=Noto Serif Italic]{Malgun Gothic}
\setsansfont{Malgun Gothic}
\setlength{\parindent}{1.2em}
\setlength{\parskip}{0.45em}
\setlength{\emergencystretch}{3em}
\begin{document}

여기서 두 문제가 본질적이다. 첫째는 \emph{주어진 표현의 환원}과 그때 나타나는 기약 성분들의 유일성 문제이다. 둘째는 주어진, 반드시 가환일 필요는 없는 체 위에서 주어진 환이 가질 수 있는 표현 전체---이를테면 기약표현만으로 한정한 전체---에 관한 문제이다.

첫째 문제는 훨씬 간단하다. 이는 표현될 환과---일반적으로 비가환인---표현체에 아무런 특별한 가정도 필요하지 않다는 점에서 이미 드러난다. 표현이 고정된 차수의 행렬들로 이루어진다는 사실이 필요한 모든 유한성 가정을 준다. 이 문제는 연산자를 가진 군의 이론으로 완전히 다룰 수 있으므로, 이제 그 이론을 간단히 살펴보겠다.

\end{document}
